\chapter{Phân tích và thiết kế hệ thống Fitty}

\section{Phân tích yêu cầu hệ thống}

\subsection{Đối tượng sử dụng}

Đối tượng sử dụng chính của hệ thống là người dùng cá nhân có nhu cầu tập luyện và theo dõi sức khỏe trên thiết bị di động. Nhóm này bao gồm người mới bắt đầu cần được hướng dẫn rõ ràng, người dùng muốn duy trì kế hoạch tập đều đặn và người dùng quan tâm tới theo dõi tiến độ theo ngày hoặc theo tuần. Ngoài ra, từ góc nhìn vận hành hệ thống, người phát triển hoặc quản trị nội dung còn đóng vai trò quản lý dữ liệu động trên Firebase như nội dung onboarding, danh mục nhóm cơ, mẫu kế hoạch khởi đầu và các cấu hình hành vi trong ứng dụng.

\subsection{Yêu cầu chức năng}

Từ việc phân tích bài toán và mục tiêu của hệ thống, các yêu cầu chức năng cốt lõi của Fitty được xác định như sau.

\begin{table}[H]
\centering
\caption{Nhóm yêu cầu chức năng chính của hệ thống Fitty}
\begin{tabular}{|p{1.5cm}|p{4.3cm}|p{7.2cm}|}
\hline
\textbf{Mã} & \textbf{Nhóm chức năng} & \textbf{Mô tả} \\ \hline
F1 & Quản lý truy cập & Cho phép người dùng đăng ký, đăng nhập bằng email, đăng nhập Google, dùng thử ở chế độ khách, khôi phục mật khẩu và đăng xuất. \\ \hline
F2 & Xác định trạng thái khởi động & Khi mở ứng dụng, hệ thống phải kiểm tra trạng thái đăng nhập và trạng thái onboarding để điều hướng về đúng màn hình ban đầu. \\ \hline
F3 & Thu thập hồ sơ ban đầu & Hệ thống phải cho phép người dùng khai báo mục tiêu, thể trạng, lịch tập, thiết bị, thói quen dinh dưỡng và nhắc việc thông qua onboarding. \\ \hline
F4 & Sinh kế hoạch tập khởi đầu & Sau onboarding, hệ thống phải tạo nội dung xem trước kế hoạch và hình thành kế hoạch tập khởi đầu cùng lịch buổi tập theo hồ sơ người dùng. \\ \hline
F5 & Quản lý nội dung bài tập & Hệ thống phải cho phép người dùng xem danh mục bài tập theo nhóm cơ, tra cứu danh sách bài tập, xem chi tiết và xem video hướng dẫn. \\ \hline
F6 & Thực hiện buổi tập & Hệ thống phải hỗ trợ người dùng bắt đầu buổi tập nhanh hoặc buổi tập theo lịch, ghi nhận kết quả bài tập và hoàn thành phiên tập. \\ \hline
F7 & Theo dõi tiến độ & Hệ thống phải tổng hợp dữ liệu buổi tập, nhật ký dinh dưỡng, chuỗi ngày hoạt động và thông tin tổng hợp theo ngày để hiển thị tiến độ cho người dùng. \\ \hline
F8 & Quản lý hồ sơ & Hệ thống phải hiển thị và cho phép cập nhật hồ sơ cá nhân, ảnh đại diện, chỉ tiêu cơ bản và một số thiết lập người dùng. \\ \hline
F9 & Thông báo & Hệ thống phải hỗ trợ thông báo cục bộ cho nhắc việc và thông báo đẩy để tăng tương tác. \\ \hline
F10 & Đồng bộ nội dung động & Hệ thống phải đọc cấu hình động từ kho dữ liệu nội dung và đồng bộ dữ liệu mô tả bài tập theo mô hình ngoại tuyến ưu tiên. \\ \hline
\end{tabular}
\end{table}

\subsection{Yêu cầu phi chức năng}

Ngoài các chức năng nghiệp vụ, hệ thống cần thỏa mãn một số yêu cầu phi chức năng quan trọng. Các yêu cầu này cũng bám theo định hướng chung của ứng dụng Android hiện đại, nơi kiến trúc, dữ liệu cục bộ và tác vụ nền cần được tổ chức thành các lớp có trách nhiệm rõ ràng \cite{android_jetpack,android_room,android_workmanager}.

\begin{itemize}
    \item \textbf{Tính khả dụng}: người dùng có thể tiếp cận các nội dung bài tập chính ngay cả khi mạng không ổn định nhờ dữ liệu cục bộ.
    \item \textbf{Tính nhất quán}: điều hướng ứng dụng phải thống nhất giữa các màn hình trước và sau khi xác thực.
    \item \textbf{Khả năng mở rộng}: kiến trúc phải cho phép bổ sung thêm nguồn dữ liệu, luồng xử lý hoặc phân hệ mới mà không phá vỡ cấu trúc hiện tại.
    \item \textbf{Khả năng bảo trì}: hệ thống cần được tách lớp rõ ràng để thuận tiện khi sửa lỗi và nâng cấp.
    \item \textbf{Hiệu năng}: dữ liệu thường dùng như metadata bài tập và thông tin người dùng cần được truy cập nhanh, hạn chế phụ thuộc tuyệt đối vào mạng.
    \item \textbf{Tính an toàn dữ liệu}: dữ liệu người dùng và phiên đăng nhập phải được quản lý có phạm vi rõ ràng giữa các tài khoản.
\end{itemize}

\section{Phân tích các phân hệ chính}

\subsection{Phân hệ truy cập và khởi động}

Phân hệ này bao gồm các màn hình Splash, Welcome, Sign In, Sign Up, Forgot Password và logic xác định trạng thái khởi động. Từ góc nhìn nghiệp vụ, đây là điểm vào của toàn bộ hệ thống. Nếu phân hệ này hoạt động sai, người dùng có thể bị điều hướng nhầm sang onboarding hoặc luồng sử dụng chính dù chưa đủ điều kiện.

Ở mức thiết kế, luồng khởi động được tổ chức theo hướng:

\begin{itemize}
    \item Splash kiểm tra trạng thái ứng dụng;
    \item nếu chưa đăng nhập thì chuyển sang Welcome hoặc Sign In;
    \item nếu đã đăng nhập nhưng chưa onboarding xong thì chuyển sang Onboarding;
    \item nếu đã hoàn tất onboarding thì chuyển sang vùng chức năng chính.
\end{itemize}

Về mặt thiết kế, phân hệ này cần bảo đảm ba yêu cầu. Thứ nhất, quyết định điều hướng ban đầu phải tập trung tại một cơ chế duy nhất để tránh mâu thuẫn giữa nhiều màn hình. Thứ hai, thông tin về trạng thái phiên đăng nhập và trạng thái khởi tạo hồ sơ phải được đọc nhanh ngay khi mở ứng dụng. Thứ ba, luồng truy cập phải cho phép mở rộng thêm các phương thức đăng nhập hoặc điều kiện khởi động mới mà không làm thay đổi cấu trúc tổng thể.

\subsection{Phân hệ onboarding và kế hoạch tập khởi đầu}

Đây là phân hệ thu thập dữ liệu đầu vào quan trọng nhất để hệ thống cá nhân hóa trải nghiệm. Các câu trả lời onboarding được dùng để:

\begin{itemize}
    \item xây dựng hồ sơ thể trạng và mục tiêu;
    \item xác lập các chỉ tiêu theo dõi cơ bản;
    \item hình thành kế hoạch tập khởi đầu và lịch tập phù hợp;
    \item tạo nền tảng cho việc hiển thị nội dung phù hợp trong các phân hệ chính.
\end{itemize}

Từ góc nhìn thiết kế, bài toán ở đây không chỉ là lưu các câu trả lời riêng lẻ mà là chuẩn hóa chúng thành một hồ sơ đầu vào nhất quán. Hồ sơ này cần đủ đơn giản để dễ mở rộng thêm trường thông tin, nhưng cũng phải đủ chặt để trở thành cơ sở cho việc hình thành kế hoạch tập, chỉ tiêu dinh dưỡng và cơ chế nhắc việc. Vì vậy, phân hệ onboarding được xem như lớp đầu vào nghiệp vụ của toàn hệ thống.

\subsection{Phân hệ nội dung bài tập}

Phân hệ này cung cấp danh sách bài tập, phân nhóm theo mục tiêu tra cứu, đồng thời hỗ trợ xem chi tiết và nội dung hướng dẫn trực quan. Điểm đáng chú ý của phân hệ là áp dụng cách tổ chức offline-first để bảo đảm người dùng vẫn có thể truy cập nội dung cốt lõi trong điều kiện mạng thay đổi.

Về mặt thiết kế, đây là phân hệ chịu tác động trực tiếp từ chất lượng kết nối mạng và dung lượng dữ liệu media. Vì vậy, hệ thống cần tách metadata bài tập khỏi các tài nguyên media, đồng thời có cơ chế đồng bộ và bộ nhớ đệm phù hợp. Cách tiếp cận này làm giảm độ trễ khi truy cập lại nội dung đã xem, đồng thời tránh việc toàn bộ trải nghiệm bị phụ thuộc hoàn toàn vào trạng thái mạng tức thời.

\subsection{Phân hệ buổi tập và lịch tập}

Phân hệ này liên quan trực tiếp đến giá trị cốt lõi của ứng dụng. Người dùng có thể bắt đầu buổi tập theo nhu cầu tức thời hoặc theo lịch đã hình thành trước đó. Khi hoàn thành phiên tập, hệ thống không chỉ kết thúc luồng hiện tại mà còn phải cập nhật các dữ liệu liên quan đến kế hoạch tập và tiến độ theo dõi. Vì vậy, đây là một use case liên phân hệ, đòi hỏi cơ chế điều phối nghiệp vụ nhất quán.

Xét theo quan điểm thiết kế, phân hệ buổi tập phải phân biệt rõ ba thực thể: kế hoạch tập ở mức tổng thể, buổi tập theo lịch và phiên tập thực tế mà người dùng đang thực hiện. Cách tách này giúp hệ thống cập nhật trạng thái linh hoạt hơn, ví dụ một buổi tập có thể được dời lịch, bỏ qua hoặc hoàn tất mà không phá vỡ cấu trúc của toàn bộ kế hoạch.

\subsection{Phân hệ theo dõi và hồ sơ}

Phân hệ theo dõi và hồ sơ có vai trò hiển thị dữ liệu người dùng dưới dạng dễ hiểu. Các thông tin tổng hợp theo ngày, ghi nhận dinh dưỡng, chỉ số cơ thể và mức độ duy trì thói quen được kết hợp lại để tạo cảm giác tiến bộ theo thời gian. Đây là nơi người dùng đánh giá hiệu quả sử dụng hệ thống sau một giai đoạn tập luyện.

Ở góc độ thiết kế, phân hệ này đóng vai trò lớp phản hồi của hệ thống. Nếu onboarding và buổi tập tạo ra dữ liệu đầu vào và dữ liệu vận hành, thì tracking và hồ sơ là nơi những dữ liệu đó được tổng hợp thành thông tin có ý nghĩa đối với người dùng. Vì vậy, mô hình tổng hợp theo ngày, chuỗi duy trì và các bản ghi tiến độ cần được tổ chức sao cho dễ tính toán, dễ mở rộng và không gây trùng lặp dữ liệu.

\section{Thiết kế kiến trúc tổng thể}

\subsection{Kiến trúc logic}

Kiến trúc logic của hệ thống có thể mô tả như Hình~\ref{fig:seq-exercise-cache}. Trước khi quan sát hình, cần lưu ý rằng luồng này được dùng để tóm tắt cách ứng dụng đồng bộ metadata bài tập và chuẩn bị media cần thiết cho trải nghiệm ngoại tuyến. Cách tổ chức này phù hợp với định hướng phân lớp của các ứng dụng di động hiện đại, trong đó giao diện, xử lý nghiệp vụ và quản lý dữ liệu cần được tách biệt để dễ mở rộng và bảo trì \cite{android_jetpack,android_compose,android_navigation_compose,firebase_firestore}.

\begin{figure}[H]
\centering
\resizebox{\textwidth}{!}{%
\begin{tikzpicture}
    \tikzset{
      seqbox/.style={draw, fill=yellow!15, minimum width=2.7cm, minimum height=1.0cm, font=\footnotesize, align=center},
      seqline/.style={dashed},
      seqmsg/.style={->, line width=0.5pt},
      seqreturn/.style={->, dashed, line width=0.5pt}
    }

    \node at (-0.2,0) {\textbf{:Người dùng}};
    \node[seqbox] (screen) at (3,0) {Màn hình\\bài tập};
    \node[seqbox] (sync) at (7,0) {Bộ đồng bộ\\bài tập};
    \node[seqbox] (remote) at (11,0) {Nguồn dữ liệu\\từ xa};
    \node[seqbox] (media) at (15,0) {Bộ đệm\\media};
    \node[seqbox] (local) at (19,0) {Kho cục bộ};

    \foreach \x in {0,3,7,11,15,19} {\draw[seqline] (\x,-0.5) -- (\x,-9.2);}

    \draw[seqmsg] (0,-1.2) -- node[above, font=\scriptsize] {1: Mở danh mục bài tập} (3,-1.2);
    \draw[seqmsg] (3,-2.0) -- node[above, font=\scriptsize] {2: Yêu cầu đồng bộ nội dung} (7,-2.0);
    \draw[seqmsg] (7,-2.8) -- node[above, font=\scriptsize] {3: Tải metadata bài tập} (11,-2.8);
    \draw[seqreturn] (11,-3.6) -- node[below, font=\scriptsize] {4: Trả metadata và liên kết media} (7,-3.6);
    \draw[seqmsg] (7,-4.4) -- node[above, font=\scriptsize] {5: Tải trước media cần thiết} (15,-4.4);
    \draw[seqmsg] (7,-5.2) -- node[above, font=\scriptsize] {6: Chuẩn hóa và lưu dữ liệu} (19,-5.2);
    \draw[seqreturn] (19,-6.0) -- node[below, font=\scriptsize] {7: Trả dữ liệu sẵn sàng dùng offline} (3,-6.0);
    \draw[seqreturn] (3,-6.8) -- node[below, font=\scriptsize] {8: Hiển thị danh mục bài tập} (0,-6.8);

    \filldraw[fill=white] (2.9,-1.4) rectangle (3.1,-6.9);
    \filldraw[fill=white] (6.9,-2.2) rectangle (7.1,-5.4);
    \filldraw[fill=white] (10.9,-3.0) rectangle (11.1,-3.8);
    \filldraw[fill=white] (14.9,-4.6) rectangle (15.1,-5.1);
    \filldraw[fill=white] (18.9,-5.4) rectangle (19.1,-6.2);
\end{tikzpicture}
}
\caption{Biểu đồ tuần tự cho luồng đồng bộ bài tập và bộ nhớ đệm media}
\label{fig:seq-exercise-cache}
\end{figure}

Biểu đồ trên cho thấy khái quát rằng luồng đồng bộ được tách thành ba bước rõ: lấy metadata, chuẩn bị media và lưu về kho cục bộ. Nhờ cách tổ chức này, danh mục bài tập có thể được khai thác ổn định hơn ngay cả khi trạng thái mạng thay đổi.

\section{Thiết kế dữ liệu}

\subsection{Mô hình dữ liệu người dùng}

Mô hình dữ liệu người dùng là trung tâm của toàn bộ hệ thống. Ở mức khái quát, một thực thể người dùng bao gồm các nhóm thông tin:

\begin{itemize}
    \item thông tin nhận diện và trạng thái tài khoản;
    \item thông tin hồ sơ thể trạng và mục tiêu;
    \item dữ liệu thu thập trong giai đoạn khởi tạo hồ sơ;
    \item thông tin về kế hoạch tập đang hoạt động;
    \item các chỉ số thống kê và theo dõi tiến độ;
    \item các thiết lập cá nhân liên quan tới trải nghiệm sử dụng.
\end{itemize}

\subsection{Mô hình dữ liệu kế hoạch tập}

Hệ thống sử dụng ba mô hình quan trọng để quản lý kế hoạch tập:

\begin{itemize}
    \item mô hình kế hoạch tập: biểu diễn kế hoạch đang gắn với người dùng;
    \item mô hình buổi tập theo lịch: biểu diễn từng buổi tập đã được phân bổ theo thời gian;
    \item mô hình thành phần bài tập: biểu diễn các bài tập cấu thành một buổi tập cụ thể.
\end{itemize}

Thiết kế này cho phép phân biệt giữa kế hoạch cấp cao và từng buổi tập cụ thể, đồng thời hỗ trợ cập nhật trạng thái buổi tập độc lập.

\subsection{Mô hình dữ liệu phiên tập và tracking}

Ở phân hệ theo dõi, hệ thống cần mô hình hóa phiên tập đã thực hiện, dữ liệu tổng hợp theo ngày và các thống kê tiến độ. Ngoài dữ liệu tập luyện, phân hệ theo dõi còn sử dụng các nhóm mô hình:

\begin{itemize}
    \item dữ liệu bữa ăn và thành phần dinh dưỡng;
    \item dữ liệu đo lường cơ thể;
    \item dữ liệu ghi nhận hình ảnh hoặc đánh giá thể trạng;
    \item dữ liệu mục tiêu và tiến độ theo ngày.
\end{itemize}

Nhờ tách mô hình như vậy, hệ thống có thể mở rộng thêm các kiểu tracking khác mà không làm rối lớp dữ liệu người dùng.

\subsection{Thiết kế dữ liệu Firebase}

Từ yêu cầu nghiệp vụ và cấu trúc dữ liệu của hệ thống, có thể tóm tắt các nhóm collection/document chính như Bảng~\ref{tab:firebase-design}.

\begin{table}[H]
\centering
\small
\caption{Thiết kế dữ liệu mức khái quát trên Firebase}
\label{tab:firebase-design}
\begin{tabular}{|p{4.9cm}|p{7.9cm}|}
\hline
\textbf{Đường dẫn dữ liệu} & \textbf{Vai trò} \\ \hline
\path|users/{uid}| & Lưu hồ sơ người dùng, trạng thái khởi tạo và các thông tin tổng hợp. \\ \hline
\path|users/{uid}/plan_instances| & Lưu các kế hoạch tập của người dùng. \\ \hline
\path|users/{uid}/plan_instances/{planId}/scheduled_workouts| & Lưu từng buổi tập theo lịch thuộc một kế hoạch. \\ \hline
\path|users/{uid}/notification_tokens| & Lưu thông tin phục vụ thông báo đẩy. \\ \hline
\path|exercises| & Lưu dữ liệu mô tả bài tập dùng cho đồng bộ theo mô hình ngoại tuyến ưu tiên. \\ \hline
\path|app_content/*| & Lưu nội dung động cho các màn hình và luồng nghiệp vụ tương ứng. \\ \hline
\path|starter_plan_templates/*| & Lưu mẫu nội dung để hình thành kế hoạch tập khởi đầu. \\ \hline
\end{tabular}
\end{table}

\subsection{Thiết kế dữ liệu cục bộ}

Ở phía thiết bị, Room đóng vai trò kho dữ liệu cục bộ cho những phần cần truy cập nhanh và có khả năng dùng lại khi ngoại tuyến \cite{android_room}. Trong đó:

\begin{itemize}
    \item bảng bài tập phục vụ tra cứu nội dung;
    \item bảng sync state phục vụ chẩn đoán trạng thái đồng bộ;
    \item bảng history hỗ trợ recently viewed;
    \item bảng notification và task phục vụ trải nghiệm thường ngày trên thiết bị.
\end{itemize}

\subsection{Biểu đồ use case tổng quan}

Ngoài phần mô tả yêu cầu bằng bảng, việc thể hiện các tương tác chính giữa tác nhân và hệ thống giúp làm rõ phạm vi nghiệp vụ của Fitty. Trong các biểu đồ use case dưới đây, quan hệ $\ll$include$\gg$ được dùng cho hành vi bắt buộc được tái sử dụng, còn quan hệ $\ll$extend$\gg$ được dùng cho hành vi mở rộng xuất hiện theo ngữ cảnh cụ thể. Hình~\ref{fig:data-overview} trình bày biểu đồ use case tổng quan, trong đó nhấn mạnh các nhóm chức năng chính dành cho người dùng và nhóm quản trị nội dung.

\begin{figure}[H]
\centering
\resizebox{\textwidth}{!}{%
\begin{tikzpicture}[
    every node/.style={font=\small,align=center},
    usecase/.style={ellipse, draw=black, fill=white, minimum width=4.4cm, minimum height=1.65cm, align=center}
]
    % actors
    \draw (-0.6,10.2) circle (0.25);
    \draw (-0.6,9.95) -- (-0.6,9.15);
    \draw (-1.05,9.65) -- (-0.15,9.65);
    \draw (-0.6,9.15) -- (-1.05,8.55);
    \draw (-0.6,9.15) -- (-0.15,8.55);
    \node[font=\small\bfseries] at (-0.6,7.85) {Người dùng};

    \draw (23.8,5.2) circle (0.25);
    \draw (23.8,4.95) -- (23.8,4.15);
    \draw (23.35,4.65) -- (24.25,4.65);
    \draw (23.8,4.15) -- (23.35,3.55);
    \draw (23.8,4.15) -- (24.25,3.55);
    \node[font=\small\bfseries] at (23.8,2.85) {Quản trị\\nội dung};

    % system boundary
    \draw[thick] (1.4,0.2) rectangle (22.0,12.3);
    \node[anchor=west, font=\bfseries] at (1.8,11.45) {Phân hệ Fitty};

    % use cases
    \node[usecase] (uc-auth) at (6.4,10.0) {Đăng nhập /\\Đăng ký};
    \node[usecase] (uc-onboarding) at (13.6,10.0) {Thực hiện\\onboarding};
    \node[usecase] (uc-plan) at (6.4,7.4) {Nhận kế hoạch\\khởi đầu};
    \node[usecase] (uc-schedule) at (13.6,7.4) {Xem kế hoạch\\tập luyện};
    \node[usecase] (uc-exercise) at (6.4,4.8) {Tra cứu\\bài tập};
    \node[usecase, minimum width=5.0cm] (uc-detail) at (13.6,4.8) {Xem chi tiết và\\media hướng dẫn};
    \node[usecase] (uc-workout) at (6.4,2.5) {Thực hiện\\buổi tập};
    \node[usecase] (uc-progress) at (13.6,2.5) {Theo dõi\\tiến độ};
    \node[usecase, minimum width=5.3cm] (uc-profile) at (6.4,0.95) {Quản lý hồ sơ và\\thống kê cá nhân};
    \node[usecase, minimum width=5.0cm] (uc-notify) at (13.6,0.95) {Nhận thông báo\\nhắc nhở};
    \node[usecase, minimum width=5.0cm] (uc-content) at (19.2,7.6) {Quản trị\\nội dung động};

    % user associations
    \draw (-0.18,9.65) -- (uc-auth.west);
    \draw (-0.18,9.3) -- (uc-onboarding.west);
    \draw (-0.18,8.9) -- (uc-plan.west);
    \draw (-0.18,8.5) -- (uc-schedule.west);
    \draw (-0.18,8.05) -- (uc-exercise.west);
    \draw (-0.18,7.6) -- (uc-workout.west);
    \draw (-0.18,7.15) -- (uc-progress.west);
    \draw (-0.18,6.7) -- (uc-profile.west);
    \draw (-0.18,6.25) -- (uc-notify.west);

    % admin associations
    \draw (23.34,4.65) -- (uc-content.east);

    % internal relations
    \draw[dashed, ->, >=stealth] (uc-plan.north east) .. controls (9.4,8.4) and (11.4,8.9) .. node[above, font=\scriptsize] {\textless\textless include\textgreater\textgreater} (uc-onboarding.south west);
    \draw[dashed, ->, >=stealth] (uc-exercise.east) -- node[above, font=\scriptsize] {\textless\textless extend\textgreater\textgreater} (uc-detail.west);
    \draw[dashed, ->, >=stealth] (uc-workout.north east) .. controls (8.4,3.8) and (11.2,6.0) .. node[below, font=\scriptsize] {\textless\textless extend\textgreater\textgreater} (uc-schedule.south west);
\end{tikzpicture}
}
\caption{Biểu đồ use case tổng quan của hệ thống Fitty}
\label{fig:data-overview}
\end{figure}

Biểu đồ trên cho thấy người dùng là tác nhân trung tâm của hệ thống, tương tác với các nhóm chức năng từ truy cập, khởi tạo hồ sơ, tra cứu bài tập, thực hiện buổi tập đến theo dõi tiến độ và quản lý hồ sơ. Ở phía còn lại, tác nhân quản trị nội dung chỉ tham gia vào use case quản trị nội dung động, phản ánh đúng ranh giới trách nhiệm giữa nhóm vận hành và người dùng cuối.

Để làm rõ hơn các cụm chức năng trọng tâm, báo cáo tiếp tục tách riêng hai nhóm use case tiêu biểu. Nhóm thứ nhất liên quan tới khởi tạo hồ sơ và hình thành kế hoạch tập. Nhóm thứ hai liên quan tới thực hiện buổi tập và theo dõi tiến độ.

\begin{figure}[H]
\centering
\resizebox{\textwidth}{!}{%
\begin{tikzpicture}[
    every node/.style={font=\small,align=center},
    usecase/.style={ellipse, draw=black, fill=white, minimum width=4.1cm, minimum height=1.65cm, align=center}
]
    \draw (-0.6,5.8) circle (0.25);
    \draw (-0.6,5.55) -- (-0.6,4.75);
    \draw (-1.05,5.25) -- (-0.15,5.25);
    \draw (-0.6,4.75) -- (-1.05,4.15);
    \draw (-0.6,4.75) -- (-0.15,4.15);
    \node[font=\small\bfseries] at (-0.6,3.45) {Người dùng};

    \draw[thick] (1.4,0.2) rectangle (15.2,7.4);
    \node[anchor=west, font=\bfseries] at (1.8,6.75) {Phân hệ khởi tạo hồ sơ và kế hoạch tập};

    \node[usecase] (uc-input) at (5.4,5.4) {Khai báo mục tiêu\\và thể trạng};
    \node[usecase] (uc-onboard-detail) at (11.1,5.4) {Thực hiện\\onboarding};
    \node[usecase] (uc-plan-detail) at (5.4,2.95) {Nhận kế hoạch\\khởi đầu};
    \node[usecase, minimum width=4.4cm] (uc-preview) at (11.1,2.95) {Xem trước kế hoạch\\và lịch tập};
    \node[usecase, minimum width=4.6cm] (uc-activate) at (8.25,1.0) {Kích hoạt kế hoạch\\ban đầu};

    \draw (-0.16,5.2) -- (uc-input.west);
    \draw (-0.16,4.9) -- (uc-onboard-detail.west);
    \draw (-0.16,4.55) -- (uc-preview.west);

    \draw[dashed, ->, >=stealth] (uc-input.east) -- node[above, font=\scriptsize] {\textless\textless include\textgreater\textgreater} (uc-onboard-detail.west);
    \draw[dashed, ->, >=stealth] (uc-onboard-detail.south west) -- node[left, font=\scriptsize] {\textless\textless include\textgreater\textgreater} (uc-plan-detail.north east);
    \draw[dashed, ->, >=stealth] (uc-plan-detail.east) -- node[above, font=\scriptsize] {\textless\textless include\textgreater\textgreater} (uc-preview.west);
    \draw[dashed, ->, >=stealth] (uc-preview.south) -- node[right, font=\scriptsize] {\textless\textless extend\textgreater\textgreater} (uc-activate.north);
\end{tikzpicture}
}
\caption{Biểu đồ use case cho cụm khởi tạo hồ sơ và kế hoạch tập}
\label{fig:usecase-onboarding-plan}
\end{figure}

Biểu đồ này cho thấy trọng tâm của giai đoạn khởi tạo là chuyển từ dữ liệu đầu vào của người dùng sang một kế hoạch tập khả dụng. Về ngữ nghĩa UML, việc khai báo mục tiêu và thể trạng là hành vi bắt buộc nằm trong use case onboarding, còn kích hoạt kế hoạch là bước mở rộng chỉ diễn ra sau khi người dùng đã xem trước kế hoạch được sinh ra.

\begin{figure}[H]
\centering
\resizebox{\textwidth}{!}{%
\begin{tikzpicture}[
    every node/.style={font=\small,align=center},
    usecase/.style={ellipse, draw=black, fill=white, minimum width=4.4cm, minimum height=1.8cm, align=center}
]
    \draw (-0.8,7.1) circle (0.25);
    \draw (-0.8,6.85) -- (-0.8,6.05);
    \draw (-1.25,6.55) -- (-0.35,6.55);
    \draw (-0.8,6.05) -- (-1.25,5.45);
    \draw (-0.8,6.05) -- (-0.35,5.45);
    \node[font=\small\bfseries] at (-0.8,4.75) {Người dùng};

    \draw[thick] (1.5,0.2) rectangle (16.2,8.1);
    \node[anchor=west, font=\bfseries] at (1.9,7.4) {Phân hệ buổi tập và theo dõi tiến độ};

    \node[usecase] (uc-sched2) at (5.8,5.95) {Xem kế hoạch\\tập luyện};
    \node[usecase] (uc-workout2) at (11.8,5.95) {Thực hiện\\buổi tập};
    \node[usecase, minimum width=5.0cm] (uc-complete2) at (5.8,3.25) {Hoàn thành phiên tập\\và lưu kết quả};
    \node[usecase] (uc-track2) at (11.8,3.25) {Theo dõi tiến độ\\theo ngày};
    \node[usecase, minimum width=5.2cm] (uc-profile2) at (8.8,1.05) {Xem hồ sơ và\\thống kê cá nhân};

    \draw (-0.36,6.6) -- (uc-sched2.west);
    \draw (-0.36,6.2) -- (uc-workout2.west);
    \draw (-0.36,5.75) -- (uc-complete2.west);
    \draw (-0.36,5.3) -- (uc-profile2.west);

    \draw[dashed, ->, >=stealth] (uc-sched2.east) -- node[above, font=\scriptsize] {\textless\textless extend\textgreater\textgreater} (uc-workout2.west);
    \draw[dashed, ->, >=stealth] (uc-complete2.north east) .. controls (8.2,4.65) and (9.7,4.95) .. node[above, font=\scriptsize] {\textless\textless extend\textgreater\textgreater} (uc-workout2.south west);
    \draw[dashed, ->, >=stealth] (uc-complete2.east) -- node[above, font=\scriptsize] {\textless\textless include\textgreater\textgreater} (uc-track2.west);
    \draw[dashed, ->, >=stealth] (uc-profile2.north east) .. controls (10.7,1.7) and (11.3,2.15) .. node[right, font=\scriptsize] {\textless\textless include\textgreater\textgreater} (uc-track2.south);
\end{tikzpicture}
}
\caption{Biểu đồ use case cho cụm buổi tập và theo dõi tiến độ}
\label{fig:usecase-workout-track}
\end{figure}

Biểu đồ này làm rõ rằng buổi tập không phải là một hành vi tách rời, mà là đầu vào trực tiếp cho hệ thống theo dõi tiến độ và thống kê hồ sơ. Việc hoàn thành phiên tập được mô hình như một nhánh mở rộng của use case buổi tập, trong khi phần theo dõi tiến độ là xử lý bắt buộc để dữ liệu giữa lịch tập, phiên tập và hồ sơ thống kê luôn nhất quán.

\begin{figure}[H]
\centering
\resizebox{\textwidth}{!}{%
\begin{tikzpicture}[
    every node/.style={font=\small,align=center},
    usecase/.style={ellipse, draw=black, fill=white, minimum width=4.1cm, minimum height=1.65cm, align=center}
]
    \draw (-0.6,6.2) circle (0.25);
    \draw (-0.6,5.95) -- (-0.6,5.15);
    \draw (-1.05,5.65) -- (-0.15,5.65);
    \draw (-0.6,5.15) -- (-1.05,4.55);
    \draw (-0.6,5.15) -- (-0.15,4.55);
    \node[font=\small\bfseries] at (-0.6,3.95) {Người dùng};

    \draw (16.0,4.9) circle (0.25);
    \draw (16.0,4.65) -- (16.0,3.85);
    \draw (15.55,4.35) -- (16.45,4.35);
    \draw (16.0,3.85) -- (15.55,3.25);
    \draw (16.0,3.85) -- (16.45,3.25);
    \node[font=\small\bfseries] at (16.0,2.65) {Quản trị\\nội dung};

    \draw[thick] (1.4,0.2) rectangle (14.8,7.5);
    \node[anchor=west, font=\bfseries] at (1.8,6.8) {Phân hệ nội dung bài tập và thông báo};

    \node[usecase] (uc-catalog) at (5.3,5.5) {Tra cứu danh mục\\bài tập};
    \node[usecase, minimum width=4.5cm] (uc-detail2) at (10.9,5.5) {Xem chi tiết và\\media hướng dẫn};
    \node[usecase] (uc-sync) at (5.3,3.05) {Đồng bộ metadata\\bài tập};
    \node[usecase] (uc-remind) at (10.9,3.05) {Nhận thông báo\\nhắc nhở};
    \node[usecase, minimum width=4.8cm] (uc-content2) at (8.1,1.0) {Quản trị nội dung động\\và cấu hình nhắc việc};

    \draw (-0.15,5.55) -- (uc-catalog.west);
    \draw (-0.15,5.15) -- (uc-detail2.west);
    \draw (-0.15,4.75) -- (uc-remind.west);

    \draw (15.56,4.35) -- (uc-content2.east);
    \draw[dashed, ->, >=stealth] (uc-catalog.east) -- node[above, font=\scriptsize] {\textless\textless extend\textgreater\textgreater} (uc-detail2.west);
    \draw[dashed, ->, >=stealth] (uc-sync.north) -- node[left, font=\scriptsize] {\textless\textless include\textgreater\textgreater} (uc-catalog.south);
\end{tikzpicture}
}
\caption{Biểu đồ use case cho cụm nội dung bài tập và thông báo}
\label{fig:usecase-content-notify}
\end{figure}

Biểu đồ này nhấn mạnh mối liên hệ giữa nội dung bài tập, dữ liệu đồng bộ và cơ chế thông báo. Về mặt chuẩn hóa use case, quản trị nội dung động được giữ là một mục tiêu độc lập của tác nhân quản trị, còn hành vi xem chi tiết bài tập chỉ là nhánh mở rộng phát sinh khi người dùng đi sâu từ danh mục sang từng bài cụ thể.

\subsection{Biểu đồ hoạt động tiêu biểu}

Để làm rõ hơn cách các phân hệ vận hành ở mức luồng nghiệp vụ, phần này trình bày một số biểu đồ hoạt động tiêu biểu. Các biểu đồ được giữ ở mức khái quát, tập trung vào các quyết định điều hướng và các bước xử lý chính thay vì đi sâu vào cài đặt.

Trước hết, Hình~\ref{fig:activity-startup} mô tả luồng khởi động ứng dụng và cách hệ thống xác định màn hình ban đầu.

\begin{figure}[H]
\centering
\begin{tikzpicture}[node distance=0.8cm, every node/.style={font=\small}, >=stealth]
    \tikzstyle{startstop}=[rounded rectangle, draw, minimum width=3.1cm, minimum height=0.8cm, align=center]
    \tikzstyle{process}=[rectangle, draw, minimum width=3.6cm, minimum height=0.8cm, align=center]
    \tikzstyle{decision}=[diamond, draw, aspect=2.1, align=center, inner sep=1pt]

    \node[startstop] (start) {Mở ứng dụng};
    \node[process, below=of start] (check) {Đọc trạng thái đăng nhập};
    \node[decision, below=of check, yshift=-0.15cm] (login) {Đã đăng nhập?};
    \node[process, left=3.8cm of login] (welcome) {Đi tới Welcome / Sign In};
    \node[process, below=of login, yshift=-0.15cm] (readonboard) {Đọc trạng thái onboarding};
    \node[decision, below=of readonboard, yshift=-0.15cm] (done) {Đã hoàn tất onboarding?};
    \node[process, left=3.9cm of done] (onboard) {Đi tới Onboarding};
    \node[process, below=of done, yshift=-0.15cm] (main) {Đi tới vùng chức năng chính};
    \node[startstop, below=of main] (end) {Hoàn tất khởi động};

    \draw[->] (start) -- (check);
    \draw[->] (check) -- (login);
    \draw[->] (login) -- node[above] {Không} (welcome);
    \draw[->] (login) -- node[right] {Có} (readonboard);
    \draw[->] (readonboard) -- (done);
    \draw[->] (done) -- node[above] {Không} (onboard);
    \draw[->] (done) -- node[right] {Có} (main);
    \draw[->] (main) -- (end);
\end{tikzpicture}
\caption{Biểu đồ hoạt động của luồng khởi động ứng dụng}
\label{fig:activity-startup}
\end{figure}

Biểu đồ này cho thấy điểm mấu chốt của phân hệ truy cập là cơ chế ra quyết định sớm. Toàn bộ ứng dụng chỉ thực sự đi vào luồng chính sau khi hệ thống đã xác định rõ cả hai điều kiện: trạng thái đăng nhập và trạng thái hoàn tất onboarding.

Tiếp theo, Hình~\ref{fig:activity-onboarding-plan} mô tả luồng onboarding gắn với việc hình thành kế hoạch tập khởi đầu.

\begin{figure}[H]
\centering
\begin{tikzpicture}[node distance=0.8cm, every node/.style={font=\small}, >=stealth]
    \tikzstyle{startstop}=[rounded rectangle, draw, minimum width=3.1cm, minimum height=0.8cm, align=center]
    \tikzstyle{process}=[rectangle, draw, minimum width=4.0cm, minimum height=0.8cm, align=center]
    \tikzstyle{decision}=[diamond, draw, aspect=2.1, align=center, inner sep=1pt]

    \node[startstop] (start) {Bắt đầu onboarding};
    \node[process, below=of start] (collect) {Thu thập mục tiêu, thể trạng, lịch tập};
    \node[process, below=of collect] (normalize) {Chuẩn hóa hồ sơ đầu vào};
    \node[decision, below=of normalize, yshift=-0.15cm] (valid) {Thông tin hợp lệ?};
    \node[process, left=4.0cm of valid] (revise) {Yêu cầu người dùng bổ sung / chỉnh sửa};
    \node[process, below=of valid, yshift=-0.15cm] (plan) {Hình thành kế hoạch tập khởi đầu};
    \node[process, below=of plan] (preview) {Hiển thị nội dung xem trước};
    \node[decision, below=of preview, yshift=-0.15cm] (accept) {Người dùng đồng ý?};
    \node[process, left=4.0cm of accept] (adjust) {Điều chỉnh lại hồ sơ hoặc lựa chọn};
    \node[process, below=of accept, yshift=-0.15cm] (activate) {Kích hoạt kế hoạch và lịch tập};
    \node[startstop, below=of activate] (end) {Hoàn tất khởi tạo hồ sơ};

    \draw[->] (start) -- (collect);
    \draw[->] (collect) -- (normalize);
    \draw[->] (normalize) -- (valid);
    \draw[->] (valid) -- node[above] {Không} (revise);
    \draw[->] (valid) -- node[right] {Có} (plan);
    \draw[->] (plan) -- (preview);
    \draw[->] (preview) -- (accept);
    \draw[->] (accept) -- node[above] {Không} (adjust);
    \draw[->] (accept) -- node[right] {Có} (activate);
    \draw[->] (activate) -- (end);
\end{tikzpicture}
\caption{Biểu đồ hoạt động của luồng onboarding và kế hoạch tập khởi đầu}
\label{fig:activity-onboarding-plan}
\end{figure}

Biểu đồ trên nhấn mạnh rằng onboarding không chỉ có vai trò thu thập thông tin, mà còn là cơ chế chuyển đổi dữ liệu đầu vào thành kế hoạch tập có thể sử dụng được. Đây là lý do phân hệ này cần được xem như một khâu tiền xử lý nghiệp vụ hơn là một tập hợp màn hình hỏi đáp đơn thuần.

Cuối cùng, Hình~\ref{fig:activity-complete-workout} mô tả luồng hoàn thành buổi tập, một luồng liên kết nhiều nhóm dữ liệu khác nhau trong hệ thống.

\begin{figure}[H]
\centering
\begin{tikzpicture}[node distance=0.8cm, every node/.style={font=\small}, >=stealth]
    \tikzstyle{startstop}=[rounded rectangle, draw, minimum width=3.1cm, minimum height=0.8cm, align=center]
    \tikzstyle{process}=[rectangle, draw, minimum width=4.0cm, minimum height=0.8cm, align=center]
    \tikzstyle{decision}=[diamond, draw, aspect=2.1, align=center, inner sep=1pt]

    \node[startstop] (start) {Người dùng kết thúc buổi tập};
    \node[process, below=of start] (save) {Ghi nhận phiên tập};
    \node[decision, below=of save, yshift=-0.15cm] (scheduled) {Có buổi tập theo lịch?};
    \node[process, left=3.8cm of scheduled] (skipupdate) {Bỏ qua cập nhật lịch};
    \node[process, below=of scheduled, yshift=-0.15cm] (updateworkout) {Cập nhật trạng thái buổi tập};
    \node[process, below=of updateworkout] (summary) {Cập nhật tổng hợp theo ngày};
    \node[process, below=of summary] (streak) {Điều chỉnh chuỗi duy trì và tiến độ};
    \node[startstop, below=of streak] (end) {Hoàn tất nghiệp vụ};

    \draw[->] (start) -- (save);
    \draw[->] (save) -- (scheduled);
    \draw[->] (scheduled) -- node[above] {Không} (skipupdate);
    \draw[->] (scheduled) -- node[right] {Có} (updateworkout);
    \draw[->] (skipupdate.south) |- (summary.west);
    \draw[->] (updateworkout) -- (summary);
    \draw[->] (summary) -- (streak);
    \draw[->] (streak) -- (end);
\end{tikzpicture}
\caption{Biểu đồ hoạt động của luồng hoàn thành buổi tập}
\label{fig:activity-complete-workout}
\end{figure}

Hình này cho thấy đặc điểm quan trọng của phân hệ buổi tập: một hành động đơn lẻ của người dùng có thể dẫn đến nhiều cập nhật trên các lớp dữ liệu khác nhau. Vì vậy, việc gom luồng này vào một cơ chế điều phối nghiệp vụ riêng là cần thiết để bảo đảm tính nhất quán của dữ liệu.

\subsection{Biểu đồ tuần tự tiêu biểu}

Bên cạnh biểu đồ hoạt động, biểu đồ tuần tự giúp làm rõ sự phối hợp giữa các thành phần ở mức khái quát. Các biểu đồ dưới đây được biểu diễn theo hướng actor--boundary--control--entity để bám sát chuẩn UML hơn và phản ánh rõ ranh giới giữa giao diện, lớp điều phối nghiệp vụ và các nguồn dữ liệu chính.

Hình~\ref{fig:seq-startup-routing} mô tả thứ tự xử lý trong luồng xác định màn hình khởi động.

\begin{figure}[H]
\centering
\resizebox{\textwidth}{!}{%
\begin{tikzpicture}
    \tikzset{
      seqbox/.style={draw, fill=yellow!15, minimum width=2.7cm, minimum height=1.0cm, font=\footnotesize, align=center},
      seqline/.style={dashed},
      seqmsg/.style={->, line width=0.5pt},
      seqreturn/.style={->, dashed, line width=0.5pt}
    }

    \node at (-0.2,0) {\textbf{:Người dùng}};
    \node[seqbox] (splash) at (3,0) {Splash\\Screen};
    \node[seqbox, minimum width=3.1cm] (usecase) at (7,0) {Khối xác định\\màn hình khởi động};
    \node[seqbox, minimum width=3.1cm] (repo) at (11,0) {Nguồn trạng thái\\khởi động và phiên};
    \node[seqbox] (nav) at (15,0) {Navigation};

    \foreach \x in {0,3,7,11,15} {\draw[seqline] (\x,-0.5) -- (\x,-8.0);}

    \draw[seqmsg] (0,-1.2) -- node[above, font=\scriptsize] {1: Mở ứng dụng} (3,-1.2);
    \draw[seqmsg] (3,-2.0) -- node[above, font=\scriptsize] {2: Yêu cầu xác định màn hình khởi động} (7,-2.0);
    \draw[seqmsg] (7,-2.8) -- node[above, font=\scriptsize] {3: Lấy startup state và kiểm tra session} (11,-2.8);
    \draw[seqreturn] (11,-3.6) -- node[below, font=\scriptsize] {4: Trả trạng thái hiện tại} (7,-3.6);
    \draw[seqmsg] (7,-4.4) -- node[above, font=\scriptsize] {5: Trả đích điều hướng} (15,-4.4);
    \draw[seqreturn] (15,-5.2) -- node[below, font=\scriptsize] {6: Điều hướng tới SignIn / Onboarding / Main} (3,-5.2);

    \filldraw[fill=white] (2.9,-1.4) rectangle (3.1,-5.4);
    \filldraw[fill=white] (6.9,-2.2) rectangle (7.1,-4.6);
    \filldraw[fill=white] (10.9,-3.0) rectangle (11.1,-3.8);
    \filldraw[fill=white] (14.9,-4.6) rectangle (15.1,-5.4);
\end{tikzpicture}
}
\caption{Biểu đồ tuần tự của luồng xác định màn hình khởi động}
\label{fig:seq-startup-routing}
\end{figure}

Biểu đồ này phản ánh một nguyên tắc thiết kế quan trọng: logic quyết định điều hướng không nên phân tán ở nhiều màn hình, mà cần được gom về một cơ chế điều phối trung tâm để giảm rủi ro sai khác trạng thái.

Hình~\ref{fig:seq-onboarding-plan} mô tả luồng từ lúc hoàn tất onboarding đến khi hình thành kế hoạch tập khởi đầu.

\begin{figure}[H]
\centering
\resizebox{\textwidth}{!}{%
\begin{tikzpicture}
    \tikzset{
      seqbox/.style={draw, fill=yellow!15, minimum width=2.7cm, minimum height=1.0cm, font=\footnotesize, align=center},
      seqline/.style={dashed},
      seqmsg/.style={->, line width=0.5pt},
      seqreturn/.style={->, dashed, line width=0.5pt}
    }

    \node at (-0.2,0) {\textbf{:Người dùng}};
    \node[seqbox] (screen) at (3,0) {Onboarding\\Screen};
    \node[seqbox, minimum width=3.3cm] (builder) at (7,0) {Khối xử lý hồ sơ\\và tạo kế hoạch};
    \node[seqbox] (content) at (11,0) {Nguồn nội dung\\động};
    \node[seqbox, minimum width=3.1cm] (userplan) at (15,0) {Kho hồ sơ\\và kế hoạch};

    \foreach \x in {0,3,7,11,15} {\draw[seqline] (\x,-0.5) -- (\x,-8.0);}

    \draw[seqmsg] (0,-1.2) -- node[above, font=\scriptsize] {1: Hoàn tất trả lời onboarding} (3,-1.2);
    \draw[seqmsg] (3,-2.0) -- node[above, font=\scriptsize] {2: Lưu câu trả lời và yêu cầu tạo kế hoạch khởi đầu} (7,-2.0);
    \draw[seqmsg] (7,-2.8) -- node[above, font=\scriptsize] {3: Lấy template và nội dung preview phù hợp} (11,-2.8);
    \draw[seqreturn] (11,-3.6) -- node[below, font=\scriptsize] {4: Trả template đã chọn} (7,-3.6);
    \draw[seqmsg] (7,-4.4) -- node[above, font=\scriptsize] {5: Lưu hồ sơ, kế hoạch và lịch tập} (15,-4.4);
    \draw[seqreturn] (15,-5.2) -- node[below, font=\scriptsize] {6: Trả preview để người dùng kích hoạt} (3,-5.2);

    \filldraw[fill=white] (2.9,-1.4) rectangle (3.1,-5.4);
    \filldraw[fill=white] (6.9,-2.2) rectangle (7.1,-4.6);
    \filldraw[fill=white] (10.9,-3.0) rectangle (11.1,-3.8);
    \filldraw[fill=white] (14.9,-4.6) rectangle (15.1,-5.4);
\end{tikzpicture}
}
\caption{Biểu đồ tuần tự của luồng hình thành kế hoạch tập khởi đầu}
\label{fig:seq-onboarding-plan}
\end{figure}

Qua biểu đồ này có thể thấy rõ rằng kế hoạch tập khởi đầu được sinh ra từ sự kết hợp giữa hồ sơ người dùng và các mẫu nội dung động, thay vì được cố định cứng trong giao diện. Cách tổ chức đó giúp hệ thống dễ bảo trì và dễ điều chỉnh khi thay đổi chiến lược nội dung.

Hình~\ref{fig:seq-complete-workout} mô tả luồng phối hợp khi người dùng hoàn thành một buổi tập.

\begin{figure}[H]
\centering
\resizebox{\textwidth}{!}{%
\begin{tikzpicture}
    \tikzset{
      seqbox/.style={draw, fill=yellow!15, minimum width=2.7cm, minimum height=1.0cm, font=\footnotesize, align=center},
      seqline/.style={dashed},
      seqmsg/.style={->, line width=0.5pt},
      seqreturn/.style={->, dashed, line width=0.5pt}
    }

    \node at (-0.2,0) {\textbf{:Người dùng}};
    \node[seqbox] (screen) at (3,0) {WorkoutSession\\Screen};
    \node[seqbox, minimum width=3.1cm] (usecase) at (7,0) {Khối hoàn thành\\buổi tập};
    \node[seqbox, minimum width=3.2cm] (planrepo) at (11,0) {Kho buổi tập\\và kế hoạch};
    \node[seqbox, minimum width=3.2cm] (trackrepo) at (15,0) {Kho dữ liệu theo dõi\\và hồ sơ};

    \foreach \x in {0,3,7,11,15} {\draw[seqline] (\x,-0.5) -- (\x,-8.0);}

    \draw[seqmsg] (0,-1.2) -- node[above, font=\scriptsize] {1: Chọn hoàn thành buổi tập} (3,-1.2);
    \draw[seqmsg] (3,-2.0) -- node[above, font=\scriptsize] {2: Gửi kết quả phiên tập} (7,-2.0);
    \draw[seqmsg] (7,-2.8) -- node[above, font=\scriptsize] {3: Hoàn tất phiên tập và cập nhật lịch tập} (11,-2.8);
    \draw[seqreturn] (11,-3.6) -- node[below, font=\scriptsize] {4: Xác nhận cập nhật kế hoạch} (7,-3.6);
    \draw[seqmsg] (7,-4.4) -- node[above, font=\scriptsize] {5: Cập nhật tổng hợp theo ngày và số buổi tập} (15,-4.4);
    \draw[seqreturn] (15,-5.2) -- node[below, font=\scriptsize] {6: Trả kết quả hoàn tất phiên tập} (3,-5.2);

    \filldraw[fill=white] (2.9,-1.4) rectangle (3.1,-5.4);
    \filldraw[fill=white] (6.9,-2.2) rectangle (7.1,-4.6);
    \filldraw[fill=white] (10.9,-3.0) rectangle (11.1,-3.8);
    \filldraw[fill=white] (14.9,-4.6) rectangle (15.1,-5.4);
\end{tikzpicture}
}
\caption{Biểu đồ tuần tự của luồng hoàn thành buổi tập}
\label{fig:seq-complete-workout}
\end{figure}

Biểu đồ trên cho thấy việc hoàn thành buổi tập là một nghiệp vụ tổng hợp, vì nó tác động đồng thời đến kế hoạch tập và dữ liệu theo dõi. Điều này củng cố yêu cầu phải có một lớp điều phối nghiệp vụ rõ ràng giữa giao diện và dữ liệu.

Hình~\ref{fig:seq-exercise-sync} mô tả thêm một luồng kỹ thuật quan trọng là đồng bộ metadata bài tập và chuẩn bị dữ liệu ngoại tuyến.

\begin{figure}[H]
\centering
\resizebox{\textwidth}{!}{%
\begin{tikzpicture}
    \tikzset{
      seqbox/.style={draw, fill=yellow!15, minimum width=2.7cm, minimum height=1.0cm, font=\footnotesize, align=center},
      seqline/.style={dashed},
      seqmsg/.style={->, line width=0.5pt},
      seqreturn/.style={->, dashed, line width=0.5pt}
    }

    \node at (-0.2,0) {\textbf{:Người dùng}};
    \node[seqbox] (screen) at (3,0) {ExerciseList\\Screen};
    \node[seqbox, minimum width=3.1cm] (repo) at (7,0) {Khối đồng bộ\\bài tập};
    \node[seqbox] (firestore) at (11,0) {Nguồn dữ liệu bài tập\\từ xa};
    \node[seqbox, minimum width=3.0cm] (cache) at (15,0) {Kho cục bộ\\và bộ đệm media};

    \foreach \x in {0,3,7,11,15} {\draw[seqline] (\x,-0.5) -- (\x,-8.0);}

    \draw[seqmsg] (0,-1.2) -- node[above, font=\scriptsize] {1: Mở danh mục bài tập} (3,-1.2);
    \draw[seqmsg] (3,-2.0) -- node[above, font=\scriptsize] {2: Yêu cầu đồng bộ metadata} (7,-2.0);
    \draw[seqmsg] (7,-2.8) -- node[above, font=\scriptsize] {3: Đọc dữ liệu mô tả bài tập từ xa} (11,-2.8);
    \draw[seqreturn] (11,-3.6) -- node[below, font=\scriptsize] {4: Trả dữ liệu thô} (7,-3.6);
    \draw[seqmsg] (7,-4.4) -- node[above, font=\scriptsize] {5: Lưu bộ đệm media và ghi dữ liệu cục bộ} (15,-4.4);
    \draw[seqreturn] (15,-5.2) -- node[below, font=\scriptsize] {6: Trả dữ liệu sẵn sàng dùng ngoại tuyến} (3,-5.2);
    \draw[seqreturn] (3,-6.0) -- node[below, font=\scriptsize] {7: Hiển thị danh sách bài tập} (0,-6.0);

    \filldraw[fill=white] (2.9,-1.4) rectangle (3.1,-6.2);
    \filldraw[fill=white] (6.9,-2.2) rectangle (7.1,-4.6);
    \filldraw[fill=white] (10.9,-3.0) rectangle (11.1,-3.8);
    \filldraw[fill=white] (14.9,-4.6) rectangle (15.1,-5.4);
\end{tikzpicture}
}
\caption{Biểu đồ tuần tự của luồng đồng bộ bài tập và dữ liệu ngoại tuyến}
\label{fig:seq-exercise-sync}
\end{figure}

Biểu đồ này bổ sung góc nhìn về phần hạ tầng dữ liệu của hệ thống. Khác với các luồng thuần nghiệp vụ, luồng đồng bộ bài tập cho thấy vai trò của cơ chế trung gian trong việc lấy dữ liệu từ xa, chuẩn hóa và lưu lại để phục vụ truy cập ổn định trên thiết bị.

\section{Thiết kế cho một số cơ chế trọng tâm}

\subsection{Hình thành kế hoạch tập khởi đầu}

Từ hồ sơ chuẩn hóa của người dùng, hệ thống:

\begin{enumerate}
    \item chọn mẫu kế hoạch phù hợp;
    \item ánh xạ thông tin hồ sơ vào mẫu nội dung;
    \item hình thành kế hoạch tập cho người dùng;
    \item tạo lịch tập tương ứng với các ngày tập;
    \item xây dựng nội dung xem trước trước khi kích hoạt kế hoạch.
\end{enumerate}

Giải pháp này có ưu điểm là phần hình thành kế hoạch được cô lập tương đối tốt và có thể thay đổi theo nội dung động, thay vì phụ thuộc trực tiếp vào giao diện.

\subsection{Đồng bộ bài tập}

Cơ chế đồng bộ bài tập được thiết kế với các bước:

\begin{enumerate}
    \item kiểm tra điều kiện đồng bộ;
    \item ghi nhận trạng thái đang đồng bộ;
    \item lấy dữ liệu bài tập từ nguồn từ xa;
    \item kiểm tra và chuẩn hóa dữ liệu;
    \item hợp nhất dữ liệu hợp lệ với dữ liệu cục bộ hiện có;
    \item lưu tạm media xem trước nếu cần thiết;
    \item cập nhật kho dữ liệu cục bộ và trạng thái đồng bộ cuối cùng.
\end{enumerate}

Thiết kế này có lợi thế là giao diện có thể phản ánh trạng thái dữ liệu một cách rõ ràng, chẳng hạn như đã đồng bộ thành công, đang dùng dữ liệu cục bộ hoặc cần đồng bộ lại.

\subsection{Hoàn thành buổi tập}

Use case hoàn thành buổi tập là ví dụ điển hình của xử lý giao dịch nghiệp vụ ở mức ứng dụng. Trình tự thiết kế được giữ theo thứ tự:

\begin{enumerate}
    \item kết thúc phiên tập hiện tại;
    \item cập nhật trạng thái của buổi tập theo lịch nếu có;
    \item cập nhật dữ liệu tổng hợp tiến độ trong ngày;
    \item điều chỉnh các chỉ số thống kê liên quan.
\end{enumerate}

Việc gom luồng này vào một cơ chế điều phối riêng giúp tránh việc phân tán logic nghiệp vụ ở nhiều thành phần giao diện khác nhau.

\section{Tổng kết chương}

Chương này đã trình bày việc phân tích yêu cầu và thiết kế hệ thống Fitty trên các góc độ chính: yêu cầu chức năng, yêu cầu phi chức năng, kiến trúc tổng thể, biểu đồ use case, biểu đồ tuần tự và thiết kế dữ liệu. Trọng tâm của chương là làm rõ cách hệ thống giải quyết bài toán ở mức kiến trúc và mô hình, trong đó nổi bật là cách tổ chức dữ liệu cho hồ sơ và kế hoạch tập, mô hình ngoại tuyến ưu tiên cho nội dung bài tập, cơ chế hình thành kế hoạch tập khởi đầu và cách điều phối các luồng nghiệp vụ chính. Trên cơ sở thiết kế đó, chương tiếp theo sẽ trình bày phần cài đặt, thực nghiệm, kiểm thử và đánh giá kết quả của hệ thống.








